\documentclass[12pt]{article}
\usepackage{geometry}                % See geometry.pdf to learn the layout options. There are lots.
\geometry{letterpaper}                   % ... or a4paper or a5paper or ... 
\usepackage{graphicx}
\usepackage{amssymb}
\usepackage{amsthm}
\usepackage{epstopdf}
\usepackage[utf8]{inputenc}
\usepackage[usenames,dvipsnames]{color}
\usepackage[table]{xcolor}
\usepackage{hyperref}
\DeclareGraphicsRule{.tif}{png}{.png}{`convert #1 `dirname #1`/`basename #1 .tif`.png}

\theoremstyle{definition}
\newtheorem{example}{Example}

\newenvironment{explanation}{%
   \setlength{\parindent}{0pt}
   \itshape
   \color{blue}
}{}

\newcommand{\projectname}{Ship4You}
\newcommand{\productname}{Ship4You}
\newcommand{\projectleader}{Alexander Hrazdera}
\newcommand{\documentstatus}{In Arbeit}
%\newcommand{\documentstatus}{Submitted}
%\newcommand{\documentstatus}{Released}
\newcommand{\version}{V. 1.0}

\begin{document}
\begin{titlepage}
\begin{flushright}
\includegraphics[scale=.5]{htlleondinglogo.png}\\
\end{flushright}

\vspace{10em}

\begin{center}
{\Huge Project Proposal} \\[3em]
{\LARGE \productname} \\[3em]
\end{center}

\begin{flushleft}
\begin{tabular}{|l|l|}
\hline
Project Name & \projectname \\ \hline
Project Leader & \projectleader \\ \hline
Document state & \documentstatus \\ \hline
Version & \version \\ \hline
\end{tabular}
\end{flushleft}

\end{titlepage}
\section*{Revisions}
\begin{tabular}{|l|l|l|}
\hline
\cellcolor[gray]{0.5}\textcolor{white}{Date} & \cellcolor[gray]{0.5}\textcolor{white}{Author} & \cellcolor[gray]{0.5}\textcolor{white}{Change} \\ \hline
September 19, 2019&A. Hrazdera/P. Schöffl&First version \\ \hline
\end{tabular}
\pagebreak

\tableofcontents
\pagebreak

\section{Einführung}
\begin{explanation}
Das Ziel dieses Projektes ist es, eine Website für Segel oder Motoryachten (Boote) zu erstellen. Der Sinn dieser Website ist es, die Boote zur Bewertung freizugeben. Somit kann jeder User der sich ein Boot mieten möchte (egal wo) ein gutes Bild über den Vermieter und den Zustand des Bootes machen. Somit kann es nicht mehr passieren das man von einem Vermieter (Charterer) über den Tisch gezogen wird. Auch kann man sich so leichter mit den Vermietern zusammenschreiben. Weil es so ein Produkt noch nicht auf dem Arbeitsmarkt gibt kann man die Idee auch leicht erweitern und vermarkten. 

\end{explanation}
\pagebreak

\section{Ausgangssituation}
\begin{explanation}
Da wir uns schon etwas mit html, css und JavaScript auskennen wird der Anfang für uns nicht so schwer werden.
Jedoch können wir Angular noch nicht und müssen uns damit erst beschäftigen. Darum werden wir uns wärend dem Programmieren und davor mit Angular auseinander setzen. Auch kennen wir uns im Team schon, weil wir schon ein anderes Projekt gemeinsam hatten. Somit können wir die kennenlern Phase auslassen.
\end{explanation}

\pagebreak

\section{General Conditions and Constraints}

\begin{example}
Das vorgeschlagene System hat mit den folgenden Einschränkungen zu kämpfen:
\begin{itemize}
\item Die Benutzerdaten sind streng vertraulich.
\item Die Benutzeroberfläche des Informationssystems muss intuitiv sein.
\item Die Anwendung ist mehrsprachig (englisch und deutsch).
\end{itemize}
\end{example}

\pagebreak

\section{Projektziele und Systemkonzepte}
\begin{example}
Die Projektziele lassen sich wie folgt zusammenfassen:
\begin{itemize}
\item Man muss nicht angemeldet sein um sich die Bewertungen zu den Booten anzusehen.
\item Anmeldung nur wenn man neue Boote hinzufügen, oder Kommentare verfassen möchte.
Auch zum Bewerten der Boote muss man sich mit Email, Passwort, Vorname, Nachname und Geburtsdatum registrieren.
\item Schnelles Fragenstellen möglich
\end{itemize}
\end{example}

\pagebreak
\section{Chancen und Risiken}
\begin{example}
Das Projekt hat folgende Möglichkeiten:
\begin{itemize}
\item Mieter bekommt für sein Geld das bestmögliche Boot
\item Risiko betrogen zu werden sinkt.
\item Man kann sich einen Überblick über das Boot machen.
\end{itemize}

Das folgende Risiko ist zu berücksichtigen.
\begin{itemize}
\item Das nicht viele Boote hinzugefügt werden
\item Das die Bewertungen verfälscht werden
\end{itemize}

\end{example}

\pagebreak
\section{Planung}
\begin{explanation}
Der erste Meilenstein wird die erstellung der Haupseite sein.
Der zweite Meilenstein wird das Anmelden der Benutzer sein.
Als drittes eine Datenbank mit den ganzen Häfen bzw. Booten.
Der vierte Meilenstein ist für die Suche dieser Häfen bzw. Boote zuständig.

Alexander Hrazdera ist für die Dokumentation und der Erstellung der Website zuständig.
Paul Schöffl ist auch für die Dokumentation und der Erstellung der Website zuständig. 

Wir benötigen 2 Programmierer mit jeweils einem eigenen Notebook. Einen Server mit einer Datenbank und eine Websitedomain wird auch benötigt.

Das Projekt wird am 30 Juni 2020 enden.
Das Projekt wird am 23 September 2019 starten. 
Answer the following questions when preparing this section:
\begin{itemize}
\item Das Projekt wird am 30 Juni 2020 enden.
\item Das Projekt wird am 23 September 2019 starten.
\item Der erste Prototyp wird ende Jänner zur Verfügung stehen.
\item Die Umsetzung wird in den nächsten Wochen beginnen.
\item Das Design der Website, der Login, die Datenbank erstellung und die Suche der Boote sind die großen Blöcke unserer Arbeit.
\item Die Arbeit wird bis zum 30 Juni beendet sein.
\item Für unsere Arbeit benötigen wir eine Datenbank und eine Website Domain.
\end{itemize}
\end{explanation}

\end{document}  
